\section{Experimental design and discovery path}
\label{sec:experimental-design}

\subsection{Objects, arithmetic, and inclusion rules}

The experiment concerns the filtration of $\theta^i f$ modulo $p^2$ for
ordinary level-one modular forms of weight $4\leq k<p-1$.  Every reported
comparison uses exact integer arithmetic.  Fourier expansions are retained
past the Sturm bound at each relevant weight, so equality of the retained
series is an exact equality of modular forms in the tested case.  No
floating-point approximation, numerical tolerance, random sampling, or
model fitting is used.

The computational record contains $5{,}572$ cells.  Of these, $5{,}434$
satisfy the ordinary condition $k<p-1$.  The main experiment
uses the $692$ ordinary, nondegenerate cells for which a $D=0$ index and its
Tate complement lie in the proved $D=2$ wedge.  These cells range over
primes $17\leq p\leq43$ and weights $8\leq k\leq22$.  A cell is included by
these arithmetic conditions alone, before its obstruction class or the
value of the candidate quadratic is inspected.

For each cell, the program performs four exact operations:
\begin{enumerate}
\item reduce $\theta^if$ modulo $p^2$ against an integral echelon basis of
      $M_{w_-}$;
\item divide the resulting canonical remainder by $p$ and reduce modulo
      $p$, producing the lower Bockstein class $\mathcal B_i$;
\item independently construct the shifted leading-block candidate
      \[
      \mathcal L_i=q_1(n,i')X^p\theta^{i-p}f,
      \qquad
      q_1(n,i')=i'^2+(k-1)i'-(n+1)(n+2);
      \]
\item compare the complete canonical remainders, not a selected coordinate
      or a short Fourier prefix.
\end{enumerate}
The program records both agreements and failures.  For each failure, it
stores the full residual
$\mathcal B_i-\mathcal L_i$ whenever the proposed identity fails.

\subsection{How the pattern was found}

We found the pattern in stages.  The first candidate reused the $D=2$
formula one Hasse step lower.  It failed in all $1{,}593$ ordinary,
nondegenerate $D=0$ cells.  This negative control showed that a new
first-order correction was needed.

Reflection to the complementary wedge then suggested replacing the earlier
quadratic $i'^2+(k-1)i'-n(n+1)$ by $q_1(n,i')$.  On the $692$ cross-wedge
cells, the resulting class agreed with the direct Bockstein in $458$ cases.
The $234$ disagreements did not scatter: they were exactly the cells on
the three affine lines
\[
 i'=n+2,
 \qquad i'+n=2p-k-1,
 \qquad i'+n=2p-k.
\]
Each line contains $78$ cells.  Away from these lines, the direct class
vanished in exactly $20$ cases, and these were exactly the cases
$q_1(n,i')=0$.  Thus the shifted conic was found only after the first
formula failed and the full residuals identified the three boundary cases.

The boundary residuals suggested further exact formulas.  On $i'=n+2$ the
residual always occupied one explicit block direction, and rational
reconstruction across primes produced the scalar
\[
 -m(k-2m-2)!,\qquad m=p-2-n.
\]
On the two triangular lines, exact elimination placed every residual in an
explicit triangular span and recovered endpoint coefficient $-(m+1)$.
These results led to the Hasse-jet and filtration arguments proved below.

\begin{figure}[t]
\centering
\resizebox{0.98\linewidth}{!}{%
\begin{tikzpicture}[
  font=\footnotesize,
  box/.style={draw=black!70, rounded corners=1pt, align=center,
              minimum height=0.85cm, inner sep=5pt},
  data/.style={box, fill=blue!10},
  found/.style={box, fill=yellow!18},
  proved/.style={box, fill=green!16},
  open/.style={box, draw=red!70!black, dashed, fill=red!4},
  flow/.style={-{Latex[length=2mm]}, thick, draw=black!65}
]
\node[data, minimum width=3.3cm] (record) at (0,0)
  {$5{,}434$ ordinary cells\\exact Sturm-safe arithmetic};
\node[data, minimum width=3.3cm, right=0.55cm of record] (wedge)
  {$692$ cross-wedge cells\\canonical Bockstein classes};
\node[found, minimum width=3.6cm, right=0.55cm of wedge] (pattern)
  {$458$ generic matches\\$234$ cells on three lines};
\draw[flow] (record) -- (wedge);
\draw[flow] (wedge) -- (pattern);

\node[proved, minimum width=4.35cm, below=0.75cm of wedge, xshift=-1.25cm]
  (theory) {proved mechanism\\first Hasse jet, perfect pairing,\\boundary-direction separation};
\node[open, minimum width=4.35cm, right=0.5cm of theory]
  (conjecture) {precise remaining target\\factorial block reduction and\\two endpoint coefficients};
\draw[flow] (pattern.south) -- ++(0,-0.28) -| (theory.north);
\draw[flow] (pattern.south) -- ++(0,-0.28) -| (conjecture.north);
\draw[flow] (theory) -- (conjecture);

\node[found, minimum width=4.35cm, below=0.65cm of conjecture]
  (tests) {$2{,}324$ third-line tests\\plus $210{,}176$ index audits\\no recorded failure};
\draw[flow] (conjecture) -- (tests);
\end{tikzpicture}
}
\caption{The experiment-to-theorem path.  Exact computation first isolated
the generic quadratic and its three boundary lines.  The green box records
the uniform results proved here.  The dashed box is the remaining symbolic
identity; finite tests support it but do not prove it.}
\label{fig:experimental-workflow}
\end{figure}



\subsection{Validation and negative controls}

Table~\ref{tab:experimental-audit} separates discovery calculations,
blind extensions, independent form tests, and symbolic-index audits.  A
zero failure count means that the stated exact identity held throughout the
specified finite range.  It does not convert the identity into a theorem
for arbitrary primes.

\begin{table}[t]
\centering
\small
\renewcommand{\arraystretch}{1.13}
\begin{tabularx}{\linewidth}{@{}>{\raggedright\arraybackslash}p{0.25\linewidth}>{\raggedright\arraybackslash}p{0.22\linewidth}>{\raggedright\arraybackslash}X@{}}
\toprule
Audit & Exact range & Outcome\\
\midrule
Ordinary-cell census & $5{,}434$ cells & Establishes the ambient finite
record and its admissible subfamilies.\\
Naive lower-block control & $1{,}593$ $D=0$ cells & The reused $D=2$
formula failed in all cells, isolating the missing projection term.\\
Cross-wedge discovery & $692$ cells & $458$ interior matches, $234$
boundary disagreements, and $20$ shifted-conic zeros.\\
Explicit first-Hasse-jet lift & The same $692$ cells & No weight,
mod-$p$, or canonical-Bockstein mismatch.\\
Third-line blind extension & $652$ cells through $p=79$ & No zero
residual, block-shape failure, or factorial-scalar mismatch.\\
Independent weight-$24$ forms & $186$ ordinary forms and $1{,}860$ tests
through $p=43$ & No zero residual, block-shape failure, or scalar mismatch.\\
Triangular endpoints & $156$ boundary cells & Endpoint rank increased by
one and the coefficient was $-(m+1)$ in every case.\\
Index and filtration audit & $210{,}176$ parameter cells through $p=101$ &
No exponent, weight, pairing, or filtration-gap failure.\\
Affine-reduction audit & $614$ factorial-range cells and $156$ recurrence
steps & No coefficient, sign, or low-point branch failure.\\
\bottomrule
\end{tabularx}
\caption{Exact audits used in the discovery and validation of the proposed
shifted-conic mechanism.  The two large validation experiments were written
as falsification tests: they count zero residuals, shape failures, and scalar
mismatches separately.}
\label{tab:experimental-audit}
\end{table}

The weight-$24$ experiment is important because the relevant cusp-form
space is two-dimensional.  It tests two generators and their projective
linear combinations rather than repeating the calculation on a single
one-dimensional family.  The third-line extension is also separate from
the original $78$ boundary cells: it searches new primes and weights after
the one-block shape and factorial scalar have been specified.

\subsection{What would falsify the conjectural mechanism}

Each computational claim has a direct failure test.  Any of the following
would disprove part of the proposed mechanism:
\begin{enumerate}
\item an interior cross-wedge cell with
      $\mathcal B_i\neq q_1(n,i')X^p\theta^{i-p}f$;
\item a cell with $q_1(n,i')=0$ and nonzero interior Bockstein, or an
      interior zero with $q_1(n,i')\neq0$;
\item a third-line residual outside the predicted one-block direction;
\item a third-line scalar different from $-m(k-2m-2)!$ or dependent on the
      chosen ordinary form;
\item a triangular residual outside the stated span or with endpoint
      coefficient different from $-(m+1)$.
\end{enumerate}
The released scripts report each failure type separately.  The tests remain
useful even though the factorial block reduction is still open.
